
\documentclass[12pt,a4paper,twoside]{book}

%% general usepackage stuff
\usepackage{times}
\usepackage{mathptmx}
\usepackage{booktabs}
\usepackage{graphicx}
\usepackage{titlesec}
\usepackage{amssymb}

\usepackage{color}

%% labels and page stuff
\renewcommand{\labelitemi}{$\diamond$}
\renewcommand{\labelitemii}{$\circ$}

\setcounter{secnumdepth}{3}

%\setlength{\oddsidemargin}{-0.5cm}
%\setlength{\evensidemargin}{-0.5cm}

\setlength{\oddsidemargin}{4.58mm}%final
\setlength{\evensidemargin}{-4.92mm}%final


\setlength{\textwidth}{15.95cm}%final
\setlength{\textheight}{23.5cm}
\setlength{\voffset}{-1.45cm} %


%% todo
\newcommand{\todo}[1]{\textcolor{red}{[todo: }#1\textcolor{red}{]}}
\newcommand{\todoh}[1]{} 
\newcommand{\done}[1]{} 

%% smallerText
\newcommand{\smallerTextSize}{10}
\newcommand{\smallerTextSkip}{12}
\newcommand{\smallerBegin}{\fontsize{\smallerTextSize}{\smallerTextSkip}\selectfont}
\newcommand{\smallerEnd}{\normalsize}
\newcommand{\smaller}[1]{\smallerBegin #1\smallerEnd}

%%reference to defintion
\newcommand{\defref}[1]{\ref{#1} on page \pageref{#1}}

%% special, short for shell
\newcommand{\shl}[1]{\sf \small #1\rm\normalsize}

%% indentation in tables
\newcommand{\tabind}[1]{\rule{#1mm}{0cm}}

%% smaller in captions
\newcommand{\captionfonts}{\smallerBegin}
                                                                                                                                                                              
\makeatletter  % Allow the use of @ in command names
\long\def\@makecaption#1#2{%
  \vskip\abovecaptionskip
  \sbox\@tempboxa{{\captionfonts #1: #2}}%
  \ifdim \wd\@tempboxa >\hsize
    {\captionfonts #1: #2\par}
  \else
    \hbox to\hsize{\hfil\box\@tempboxa\hfil}%
  \fi
  \vskip\belowcaptionskip}
\makeatother   % Cancel the effect of \makeatletter

% % clear page before new chapter
% \makeatletter
% \def\cleardoublepage{\clearpage\if@twoside \ifodd\c@page\else
% \hbox{}
% \vspace*{\fill}
% \begin{center}
% %This page intentionally contains only this sentence.
% \end{center}
% \vspace{\fill}
% \thispagestyle{empty}
% \newpage
% \if@twocolumn\hbox{}\newpage\fi\fi\fi}
% \makeatother

%% abbreviations
\usepackage{nomencl}
\let\abbrev\nomenclature
\renewcommand{\nomname}{List of Abbreviations}
\setlength{\nomlabelwidth}{.25\hsize}
\renewcommand{\nomlabel}[1]{#1 \dotfill}
\setlength{\nomitemsep}{-\parsep}
\makenomenclature
\newcommand{\Listofabbrev}{
\printnomenclature
\newpage
}

%% chapter title formatting
\titleformat{\chapter}[display]{ \raggedleft }{\fontsize{52}{63}\selectfont \bf \thechapter }{0.2cm}{\fontsize{32}{38.7}\selectfont  }[]

%% haeder formatting
% \renewcommand{\chaptermark}[1]{%
% \markboth{\chaptername
% \ \thechapter.\ #1}{}}
\renewcommand{\sectionmark}[1]{\markright{\thesection.\ #1}}
\usepackage{fancyhdr}
\pagestyle{fancy}
\fancyhf[LEH,ROH]{\thepage}
\fancyhf[REH]{\smaller{\nouppercase{\leftmark}}}
\fancyhf[LOH]{\smaller{\it \nouppercase{\rightmark}}}
\fancyhf[COF]{\rule{0.2cm}{0.0cm}}
\fancyhf[CEF]{\rule{0.2cm}{0.0cm}}
\renewcommand{\headrulewidth}{0pt}

%% macro for figures
%\usepackage{svg}
%\usepackage{amsmath}
%\newcommand{\printlabel}{}
%\newcommand{\abcdef}[1]{\tiny #1 \normalsize}

		%% arguments: graphics file, label, caption, smallcaption
\newcommand{\insertFigure}[4]{\begin{figure}[top] \smallerBegin \centering \includegraphics{#1}  \\  \caption{\label{#2}\smallerBegin #3 \footnotesize{#4}  \smallerEnd }  
\end{figure}	}
%% macro for figures with short caption
		%% arguments: graphics file, label, caption, smallcaption, shortcaption
\newcommand{\insertFigureShort}[5]{\begin{figure}[top] \smallerBegin \centering \includegraphics{#1} \label{#2} \\ \caption[#5]{\smallerBegin #3 \footnotesize{#4} \smallerEnd } 
\end{figure}	}
%\renewcommand{\topfraction}{1}

\newcommand{\spaceafterpar}{\vspace{14.48pt}}

\renewcommand{\floatpagefraction}{.75} % vorher: .5
\renewcommand{\textfraction}{.1}       % vorher: .2
\renewcommand{\topfraction}{.8}        % vorher: .7
\renewcommand{\bottomfraction}{.5} 
\setcounter{topnumber}{3}              % vorher: 2
\setcounter{bottomnumber}{2}           % vorher: 1
\setcounter{totalnumber}{5}            % vorher: 3
%from: http://www.matthiaspospiech.de/latex/vorlagen/allgemein/preambel/9/

%macros for tables
    %% arguments: columns
\newcommand{\tableBegin}[1]{\begin{table}[top] \begin{center} \smallerBegin \begin{tabular}{#1}}
		%% arguemtns: caption, label
\newcommand{\tableEnd}[2]{ \end{tabular} \smallerEnd \end{center} \caption{#1} \label{#2}	\end{table}}

%% shortcuts
\newcommand{\emit}[1]{\item \emph{#1}:}
\newcommand{\firstappear}[2]{\emph{#1} (#2) \abbrev{#2}{#1}}

\hyphenation{or-gan-izing}

%% definitions

\newtheorem{definition}{Definition}

%% examples
\newcommand{\exampleBeginText}[1]{\paragraph{#1}}
\newcommand{\exampleEnd}{\vspace{6mm}}
\newcommand{\exampleBegin}{\exampleBeginText{Example}}

%algorithms
\usepackage[english,ruled,vlined, slide, norelsize]{algorithm2e}

%references

\usepackage[numbers,sort]{natbib}
\setlength{\bibsep}{0.0pt}
% \bibliographystyle{plainnat}

%ref imgs
\usepackage{hyperref}
\usepackage{cleveref}

\usepackage{csvsimple}


% Front Matter

\newcommand{\titleLineOne}{Implentation of a Simplistic Interface for}
\newcommand{\titleLineTwo}{Big Data Workload Scheduler in Kubernetes}
\newcommand{\titleLineThree}{}
\newcommand{\documentdate}{March 29, 2022}
\newcommand{\studentname}{Lukas Schwerdtfeger}
\newcommand{\abstracttextde}{Das Datenaufkommen hat sich in den letzten Jahren vervielfältigt  und wird sich voraussichtlich weiter beschleunigen. Die Verarbeitung von Big Data erfordert hochentwickelte Systeme, die mit dem Datenfluss Schritt halten können. Spark und Flink sind oft genutzte Big-Data-Frameworks, die die Entwicklung von verteilten Datenflussanwendungen unterstützen. Spark- und Flink-Anwendungen werden in der Regel auf einem gemeinsam genutzten Cluster eingesetzt, der aus preiswerter Standardhardware besteht. Bisher wurde eine statische Verteilung von Cluster-Ressourcen verwendet, um langlaufende Batch-Anwendungen voneinander zu trennen, die aber mit potenziell ungenutzter Hardware einhergeht. Cluster-Ressourcenmanager wie YARN, Mesos oder Kubernetes ermöglichen eine feinere Verteilung gemeinsam genutzter Cluster-Ressourcen unter Verwendung von sogenannten Containern, so dass Anwendungen dynamischer skaliert werden können, wenn Ressourcen verfügbar sind. Eine feine Aufteilung schafft viele Möglichkeiten, Resourcenverteilungen für Batch-Anwendungen zu planen. Die Entwicklung eines neuen Planungsalgorithmus erfordert ein tiefes Verständnis der inneren Funktionsweise des Clusters. Die Suche nach einer optimalen Resourcenverteilung hat erhebliche Auswirkungen auf die Gesamtressourcennutzung des Clusters und die Laufzeit der Jobs. Kubernetes ist ein Ressourcenmanager, der sich heute als Industriestandard durchgesetzt hat. Seine Schnittstelle unterscheidet sich erheblich von anderen Ressourcenmanagern wie YARN. Planungsalgorithmen, die für YARN entwickelt wurden, können nicht ohne wesentliche Änderungen auf Kubernetes angewendet werden. Diese Arbeit wird eine einfache Schnittstelle zur Steuerung des Scheduling von Batch-Anwendungen auf Kubernetes entwickelt. Die Schnittstelle erweitert Kubernetes um TestBeds, die Slots in einem Cluster reservieren und sie einem externen Scheduler anbieten. Der externe Scheduler übermittelt sein Scheduling, dass die Jobs in die Warteschlange stellt und sie in der Reihenfolge ihrer zugewiesenen Slots ausführt. Die Schnittstelle wird durch die Implementierung eines Profiling-Schedulers mit Hilfe mehrerer Testbeds evaluiert.}
\newcommand{\abstracttext}{Data generation has exploded over the last few years and is not expected to slow down. Big Data processing requires sophisticated systems that can keep up with the data flow. Spark and Flink are popular Big Data frameworks that aid the development of distributed dataflow applications. Spark and Flink applications are commonly deployed on a shared cluster composed of cheap commodity hardware. Previously, static partitioning of cluster resources has been used to isolate long-running batch applications from one another, which comes with potentially unused hardware. Cluster resource managers like YARN, Mesos, and Kubernetes allow a more fine-grain partitioning of shared cluster resources, using containers, allowing applications to scale more dynamically if resources are available. A fine-granular partitioning creates many possible schedulings of batch applications. Developing a new scheduling algorithm requires a deep understanding of the cluster's inner workings. Finding optimal scheduling will significantly impact the cluster's overall resource utilization and runtime of jobs. Kubernetes is a resource manager that has become the defacto industry standard. Its interface differs considerably from other resource managers like YARN. Scheduling algorithms that have been developed for YARN can not be applied to Kubernetes without substantial changes. This work offers a simplistic interface to control the scheduling of batch applications on Kubernetes. The interface extends Kubernetes with Testbeds reserving slots in a cluster and offering them to an external scheduler. The external scheduler submits its scheduling, which will enqueue jobs and run them in order inside their dedicated slots. The interface is evaluated by implementing a profiling scheduler with the help of multiple Testbeds.}
\newcommand{\acktext}{This chapter is optional. First of all, I would like to...}
\providecommand{\tightlist}{%
  \setlength{\itemsep}{0pt}\setlength{\parskip}{0pt}}


\begin{document} 

\thispagestyle{empty}
\rule{0cm}{2cm}

\fontsize{19.83}{26.45} \selectfont 
\noindent \textbf{\titleLineOne} \vspace{0.2cm}
\\\textbf{\titleLineTwo}\vspace{0.2cm}
\\\textbf{\titleLineThree}
\normalsize

%\includegraphics[scale=0.9]{graphics/1445597046599wholeGraph3.pdf} 


\newpage
\thispagestyle{empty}
\rule{0cm}{2cm}
\newpage

\thispagestyle{empty}
\rule{0cm}{1.5cm}

\fontsize{24.44}{32} \selectfont 
\noindent \textbf{\titleLineOne} \vspace{0.2cm}
\\\textbf{\titleLineTwo}\vspace{0.2cm}
\\\textbf{\titleLineThree}
\normalsize

%\noindent ---Draft--- \vspace{-0.75cm}


\vspace{1.5cm}
 
\noindent \textbf{\studentname}

\vspace{-0.25cm}
\vspace{7cm}
%\vspace{10cm}

\noindent A thesis submitted to the 
\\\textbf{Faculty of Electrical Engineering and Computer Science
} 
\\of the 
\\\textbf{Technical University of Berlin} 
\\in partial fulfillment of the requirements for the degree 
\\\textbf{Bachelor Technische Informatik}

\vspace{-0.5cm}
 
\noindent Berlin, Germany\\ 
\documentdate\\

\vspace{-0.83cm}
%\vspace{1cm}

\noindent \includegraphics[width=30mm,angle=0]{graphics/tu_logo}

\rule{0cm}{20cm}

\noindent First supervisor: 
\noindent Prof. Dr. habil. Odej Kao, Technical University of Berlin

\noindent Second supervisor: 
\noindent Prof. Dr. rer. nat. Volker Markl, Technical University of Berlin

%\noindent Date of public examination: November 20, 2009
\thispagestyle{empty}

\rule{0cm}{10cm}

\noindent Hiermit erkl\"are ich, dass ich die vorliegende Arbeit selbstst\"andig und eigenh\"andig sowie ohne
unerlaubte fremde Hilfe und ausschließlich unter Verwendung der aufgef\"uhrten Quellen und
Hilfsmittel angefertigt habe.

Berlin, den

\vspace{1cm}
...........................................................................

Unterschrift
\vspace{1cm}


%\noindent Date of public examination: November 20, 2009
\thispagestyle{empty}

%\chapter*{Selbstst\"andigkeitserkl\"arung}

%\selbststaendigkeitserklaerung

%\thispagestyle{empty}
%\newpage
%\vspace*{3cm}
%\thispagestyle{empty}

\chapter*{Zusammenfassung}

\abstracttextde

\thispagestyle{empty}
\newpage
\vspace*{3cm}
\thispagestyle{empty}

\chapter*{Abstract}

\abstracttext

\thispagestyle{empty}
\newpage
\vspace*{3cm}
\thispagestyle{empty}

%% indices


%\thispagestyle{empty}
\pagenumbering{roman}\setcounter{page}{8}
\tableofcontents

%\thispagestyle{empty}
\newpage


\listoffigures
\listoflistings
\newpage
\thispagestyle{empty}
\newpage
\vspace*{3cm}
\thispagestyle{empty} \newpage
\pagenumbering{arabic}\setcounter{page}{1}

% Body Matter (use input to add chapters)


\chapter{Introduction}
\section{Motivation}
\input{generated/Introduction/Motivation.tex}
\section{Problem Description}
\input{generated/Introduction/Problem_Description.tex}
\section{Goal of this Thesis}
\input{generated/Introduction/Goal_of_this_Thesis.tex}
\section{Structure of this Thesis}
\input{generated/Introduction/Structure_of_this_Thesis.tex}


\chapter{Background}
\section{Distributed Dataflow Applications}
\input{generated/Background/Big_Data_Stream_Processing.tex}
\section{Scheduling}
\input{generated/Background/Scheduling.tex}
\section{Cluster Management Systems}
\input{generated/Background/Cluster_Management_Systems.tex}

\chapter{Approach}
\input{generated/Approach/init.tex}
\section{Extending Kubernetes using the Operator Pattern}
\input{generated/Approach/Extending_Kubernetes_using_the_Operator_Pattern.tex}
\section{Scheduling in Kubernetes}
\input{generated/Approach/Scheduling_in_Kubernetes.tex}


\chapter{Implementation}
\section{Designing the Interface}
\input{generated/Implementation/Designing_the_Interface.tex}
\section{Architecture}
\input{generated/Implementation/Architecture.tex}
\section{Operators}
\input{generated/Implementation/Operator.tex}

\chapter{Evaluation}
\section{Testing}
\input{generated/Evaluation/Testing.tex}
\section{Example Scheduling Algorithm} \label{sec:exampleSchedulingAlgorithm}
\input{generated/Evaluation/Example_Scheduling_Algorithm.tex}
\section{Limitations}
\input{generated/Evaluation/Limitations.tex}

\chapter{State of the Art}
\section{Volcano}
\input{generated/State_of_the_Art/Volcano.tex}
\section{Flux}
\input{generated/State_of_the_Art/Flux.tex}

\chapter{Conclusion and Future Work}
\section{Conclusion}
\input{generated/Conclusion/Conclusion.tex}
\section{Future Work}
\input{generated/Conclusion/Future_Work.tex}


% Back Matter (use input to add appendices, if really needed)

\appendix
\chapter{Repositories}
\input{generated/Appendix/Repository.tex}

\newpage
\smallerBegin
\pagenumbering{roman}
\bibliographystyle{plain}
\nocite{*}
\bibliography{all_publications}

\smallerEnd

%% end of document

\end{document}